% \iffalse
%<*driver>
\ProvidesFile{seqenseigne-core.dtx}%
[2025/10/24 v1.0-dev seqenseigne – Module core (coeur fonctionnel)]
%</driver>
% \fi
%
% =============================================================
% seqenseigne-core.dtx
% Module : CORE
% Rôle   : coeur logique du package seqenseigne :
%          - version et meta
%          - options du package
%          - chargement des dépendances LaTeX communes
%          - macros structurelles "publiques" \seq...
%          - macros internes \seq@...
%          - environnements structurants (ex: série d'exercices,
%            blocs de séquence, etc.)
%
% Convention de nommage :
%  - commandes publiques    \seqTruc
%  - commandes internes     \seq@Truc
%  - couleurs               \seqCouleurTruc (→ plutôt dans theme, mais
%                           les variables globales peuvent être déclarées ici)
%
% Toute macro ou environnement utilisé activement dans les livrets /
% cartes mentales / plans de travail va ici, sauf si :
%  - c'est purement visuel          → seqenseigne-theme.dtx
%  - c'est du chargement de données → seqenseigne-data.dtx
%  - c'est de l'ancien code         → seqenseigne-legacy.dtx
%
% =============================================================
% Structure des documents de référence :
% ------------------------
% Livret de séquence
% ------------------------
% Exercice
% ------------------------
% Connaissance
% ------------------------
% Savoir-faire
% ------------------------
% Livret de plans de travail
% ------------------------
% Fiches de résumé à compléter
% ------------------------
% Bilan par compétence
% ------------------------
% Evaluation notée
% ------------------------
% Suivi autonome de progression
% =============================================================
% Structure des documents de suivi de classe :
% ------------------------
% Suivi annuel par élève
% ------------------------
% Suivi annuel par séquence
% =============================================================
%
%<*package>
\NeedsTeXFormat{LaTeX2e}%
\ProvidesPackage{seqenseigne-core}%
[2025/10/24 v1.0-dev seqenseigne core module]

%% ------------------------------------------------------------------
%% ---- Chargement des extentions  ----------------------------------
%% ------------------------------------------------------------------
% --- Saisie en UTF8
\RequirePackage[utf8]{inputenc}%
% --- Police en Latin Modern 1 avec les caractères accentués et 
% --- résolution des pb de pixelisation avec lmodern
\RequirePackage[T1]{fontenc}%
\RequirePackage{lmodern}%
% --- Traduction en français et commandes spécifiques
\RequirePackage[french]{babel}%
% --- Réglage de la taille de la page
\RequirePackage{geometry}%
% --- Symboles mathématiques
\RequirePackage{amsmath}%
\RequirePackage{amssymb}% \leqslant, \geqslant, \mathbb, etc. utilisés dans les exercices
% amsfonts supprimé : chargé automatiquement par amssymb
% --- Calcul des pages (numéros, entêtes, pieds de page)
\RequirePackage{lastpage}%
\RequirePackage{fancyhdr}%
% --- plusieurs colonnes par pages, plusieurs lignes regroupées dans un tableau
\RequirePackage{multicol}%
\RequirePackage{multirow}%
% --- Gestion de tableaux
% tabularx supprimé : 0 usage dans le paquet (les documents peuvent le charger si besoin)
\RequirePackage{tabularray} \UseTblrLibrary{booktabs}%
\RequirePackage{xltabular}%
\RequirePackage{booktabs}%
% --- Gestion des couleurs
\RequirePackage{colortbl}%
\RequirePackage{xcolor}%
% --- Gestion des commandes et des environnements
\RequirePackage{etoolbox}%
\RequirePackage{xkeyval}%
% --- Mise en forme conditionnelle, boucles
\RequirePackage{xifthen}%
\RequirePackage{forloop}%
% --- Graphiques
\RequirePackage{tikz}%
\usetikzlibrary{babel}%  compatibilite tikz+babel
% tkz-euclide, tkz-tab, tkz-base : charges via option [geometrie]
% (voir \DeclareOption{geometrie} ci-dessous)
% kvoptions supprimé : 0 usage dans le paquet
% --- Liens hypertexte
\RequirePackage{hyperref}%
% scratch3 chargé via option [scratch] (voir DeclareOption ci-dessous)
% --- Symbole euro
\RequirePackage{eurosym}%
\RequirePackage{environ}%
\RequirePackage{xintexpr}%
% ---- Format de date
\RequirePackage{datetime}%

%% ------------------------------------------------------------------
%% ---- Traitement des options --------------------------------------
%% ------------------------------------------------------------------
\newif\ifdocstd\docstdtrue
\DeclareOption{ltxdoc}{\docstdfalse}%
%% Option geometrie : charge tkz-euclide, tkz-base, tkz-tab
%% Utilisation : \usepackage[geometrie]{seqenseigne}%
%% A ajouter dans les livrets/exercices de géométrie (S10-S13 typiquement)
\newif\ifseq@geometrie\seq@geometriefalse
\DeclareOption{geometrie}{\seq@geometrietrue}%
%% Option scratch : charge scratch3 (blocs Scratch, algorithmique)
%% Utilisation : \usepackage[scratch]{seqenseigne}%
%% Nécessaire uniquement pour les séquences d'algorithmique
\newif\ifseq@scratch\seq@scratchfalse
\DeclareOption{scratch}{\seq@scratchtrue}%
\ProcessOptions\relax
\ifseq@geometrie
    \RequirePackage{tkz-base}%
    \RequirePackage{tkz-euclide}%
    \RequirePackage{tkz-tab}%
    \usetikzlibrary{intersections,mindmap}%
\fi
\ifseq@scratch
    \RequirePackage{scratch3}%
\fi

% -------------------------------------------------------------
% Options globales du package
% (si tu avais \DeclareOption etc. dans ton dtx d'origine,
% copie-les dans ce bloc)
% -------------------------------------------------------------
% \DeclareOption{...}{...}%
% \ProcessOptions\relax

%% ------------------------------------------------------------------
%% ---- Variables globales ------------------------------------------
%% ------------------------------------------------------------------
% Ces commandes peuvent être redéfinies dans macro_config.tex
\providecommand{\nomCollege}{Établissement}%
\providecommand{\codeNiveau}{N00}%
\providecommand{\numNiveau}{numéro du niveau}%
\providecommand{\nomNiveau}{nom du niveau}%
\providecommand{\codeNiveauPrec}{N00}%
\providecommand{\numNiveauPrec}{numéro du niveau précédent}%
\providecommand{\nomNiveauPrec}{niveau précédent}%
\providecommand{\nomClasse}{Classe}%
\providecommand{\annee}{2024/2025}%
\providecommand{\dirPrefix}{.}%
\providecommand{\exoDir}{../Exos}%
\providecommand{\codeObj}{Objectif 00}%
\providecommand{\seqCodeObj}{Objectif 00}%
%
%% ------------------------------------------------------------------
%% ---- Variables internes : niveau et séquence courants -----------
%% ------------------------------------------------------------------
% \begin{macro}{\seq@code@niveau}%
% \begin{macro}{\seq@code@sequence}%
% Variables internes (\gdef) : niveau et séquence en cours de traitement.
% Ne pas utiliser directement — passer par \seqCodeNiveau et \seqCodeSequence.
% \begin{macrocode}%
\gdef\seq@code@niveau{N00}%
\gdef\seq@code@sequence{S00}%
% \end{macrocode}%
% \end{macro}%
% \end{macro}%
%
% \begin{macro}{\seqCodeNiveau}%
% \begin{macro}{\seqCodeSequence}%
% Accesseurs publics en lecture (expansibles).
% \begin{macrocode}%
\newcommand{\seqCodeNiveau}{\seq@code@niveau}%
\newcommand{\seqCodeSequence}{\seq@code@sequence}%
% \end{macrocode}%
% \end{macro}%
% \end{macro}%
%
% \begin{macro}{\seqSetCodeNiveau}%
% \IndexMacro{\seqSetCodeNiveau}%
% \begin{macro}{\seqSetCodeSequence}%
% \IndexMacro{\seqSetCodeSequence}%
% Mutateurs publics.
% #1 : nouveau code niveau (ex. \verb|N10|) ou séquence (ex. \verb|S03|).
% \begin{macrocode}%
\newcommand{\seqSetCodeNiveau}[1]{\gdef\seq@code@niveau{#1}}%
\newcommand{\seqSetCodeSequence}[1]{\gdef\seq@code@sequence{#1}}%
% \end{macrocode}%
% \end{macro}%
% \end{macro}%
%
%% ------------------------------------------------------------------
%% ---- Gestion des chemins de fichiers ----------------------------
%% ------------------------------------------------------------------
% \begin{macro}{\seqDefChemin}%
% \IndexMacro{\seqDefChemin}%
% Définit les chemins pour \input et \graphicspath
% #1 : liste de chemins (format : {{chemin1}{chemin2}...})
% \begin{macrocode}%
\newcommand{\seqDefChemin}[1]{%
    \def\input@path{#1}%
    \graphicspath{#1}%
}%
% \end{macrocode}%
% \end{macro}%
%
% \begin{macro}{\seqDefCheminImages}%
% \IndexMacro{\seqDefCheminImages}%
% Définit un dossier centralisé pour les images, en complément du
% \cs{graphicspath} déjà posé par \cs{seqDefChemin}.
%
% Conçu pour le scénario de la version 0.7 où toutes les images sont
% rassemblées dans un dossier plat (typiquement
% \texttt{appli/data/images/}) et nommées d'après le \texttt{.tex} qui
% les utilise. Cette macro \emph{ajoute} le chemin sans écraser ceux
% déjà déclarés par \cs{seqDefChemin}, ce qui permet aux deux
% conventions de cohabiter pendant la transition.
%
% Variable interne : \cs{seq@chemin@images} (utilisable par d'autres
% macros du paquet pour résoudre une image par son nom).
%
% \#1 : chemin absolu ou relatif vers le dossier d'images
%       (sans accolades doubles, on ajoute nous-mêmes le wrapping).
%
% Exemple : \verb|\seqDefCheminImages{D:/Enseignement/seqenseigne/appli/data/images/}|
% \begin{macrocode}%
\gdef\seq@chemin@images{}%
\newcommand{\seqDefCheminImages}[1]{%
    \gdef\seq@chemin@images{#1}%
    \g@addto@macro\Ginput@path{{#1}}%
}%
% \end{macrocode}%
% \end{macro}%
%
% \begin{macro}{\seqCheminImages}%
% \IndexMacro{\seqCheminImages}%
% Accesseur public en lecture sur \cs{seq@chemin@images}. Renvoie la
% chaîne vide si aucun chemin n'a été défini par
% \cs{seqDefCheminImages}.
% \begin{macrocode}%
\newcommand{\seqCheminImages}{\seq@chemin@images}%
% \end{macrocode}%
% \end{macro}%
%
% =============================================================
% Macros structurelles / "API publique"
% =============================================================
% Ici : toutes les commandes \seq... qui décrivent la structure
% pédagogique d'une séquence, sans faire appel directement au style.
%
% EXEMPLES À COLLER DE TON DTX :
%   - \seqTitreSequence
%   - \seqNiveauCourant / \seqSequenceCourante (ou équivalent)
%   - les macros qui construisent les en-têtes de livret
%   - gestion d'encarts de type "Objectifs", "Pré-requis" si la mise
%     en forme n'est pas directement graphique
%
% Chaque macro doit être documentée puis codée dans un bloc macrocode.
%
% \begin{macro}{\seqNomDeSequence}%
% \IndexMacro{\seqNomDeSequence}%
% Description pédagogique...
% \begin{macrocode}%
% \newcommand{\seqNomDeSequence}{...}%
% \end{macrocode}%
% \end{macro}%
%
% =============================================================


%% ------------------------------------------------------------------
%% ---- Plan de travail ---------------------------------------------
%% ------------------------------------------------------------------

% -----------------------------------------------
% \seqBoiteTitrePlan
%  - Titre du plan de travail de la séquence 
%  - Utilise \seqCodeSequence, \annee \nomClasse et \nomCollege
% -----------------------------------------------
%    \begin{macro}{\seqBoiteTitrePlan}%
%% ---- Options 
\define@cmdkey[seq]{boiteTitrePlan}{afficheSeq}[oui]{}%
\define@cmdkey[seq]{boiteTitrePlan}{affichePlan}[oui]{}%
\define@cmdkey[seq]{boiteTitrePlan}{titre}{}%
\newcommand{\seqBoiteTitrePlan}[1][]{%
	\setkeys[seq]{boiteTitrePlan}{afficheSeq,affichePlan,#1}%
	\begin{boiteTitreGen}{\ifthenelse{\equal{\cmdseq@boiteTitrePlan@afficheSeq}{oui}}{\large \textbf{Séquence \seqCodeSequence}}{}}%
		\begin{minipage}[c][50pt][c]{\linewidth}%
			\vfill { \small \textit{\ifthenelse{\equal{\cmdseq@boiteTitrePlan@affichePlan}{oui}}{Plan de Travail}{}} \hfill \textit{Année \annee} }%
			\vfill \centering {\Large \textbf{\cmdseq@boiteTitrePlan@titre}}%
			\vfill { \small \hfill \textit{\nomCollege\;-\;\nomClasse} }%
		\end{minipage}%
	\end{boiteTitreGen}%
}%
%    \end{macro}%



% -----------------------------------------------
% \seqPlanObjectif
%  - Bloc repères temporels
%  - Bloc niveaux atteints
%  - Utilise \seq@code@niveau et \seq@code@sequence (variables globales)
% Paramètres:
%  #1 : code objectif (Ex: 02)
%  #2 : Connaissances (Ex: \og \seqConnaissanceGetNom{\seq@code@niveau}{\seq@code@sequence}{03} \fg{})
%  #3 : numéros des exos série F (Ex: 1,2)
%  #4 : numéros des exos série A (Ex: 1,2)
%  #5 : numéros des exos série E (Ex: 1,2)
% -----------------------------------------------
\newcommand{\seqPlanObjectif}[5]{%
	\begin{boitePaleNoBreak}{\begin{minipage}{0.6\textwidth} \textbf{#1}: \seqObjectifGetNom{#1} \end{minipage}}%
		\begin{seqColItem}[nbCols=1,label=\ding{109}]
			\item Notions: #2
		\end{seqColItem}%
		\begin{itemize*}[label=\qquad\ding{109}]
			\item Série fondamentale: #3 \item Série avancée: #4	\item Série exploration: #5
		\end{itemize*}%
	\end{boitePaleNoBreak}%
}%


%% ------------------------------------------------------------------
%% ---- Fiche PPRE (Plan Personnalisé de Réussite Éducative) --------
%% ------------------------------------------------------------------
%
% Une fiche PPRE accompagne un élève bénéficiant d'un plan personnalisé
% de réussite. Pour chaque objectif d'une séquence, elle rappelle les
% exercices du livret (numéros F / A / E, identiques à \seqPlanObjectif)
% et ajoute, pour chaque série, une aide spécifique à cet élève :
% indice de démarrage, guide de rédaction, mémo-formule, schéma, etc.
%
% La zone d'aide est entièrement libre (contenu LaTeX quelconque).
% Si elle est vide, la boîte d'aide n'est pas affichée.
%
% UNE PAGE PAR OBJECTIF.
%
% -----------------------------------------------------------------------
% \seqPPRE — macro principale
% -----------------------------------------------------------------------
% Options (toutes passées dans l'argument optionnel clé=valeur) :
%
%   niveau     : code niveau            (ex. N11)         [obligatoire]
%   sequence   : code séquence          (ex. S03)         [obligatoire]
%   objectif   : code objectif          (ex. 02)          [obligatoire]
%   exosF      : numéros exos série F   (ex. {1,2,3})     [obligatoire]
%   exosA      : numéros exos série A   (ex. {4,5})       [obligatoire]
%   exosE      : numéros exos série E   (ex. {6})         [obligatoire]
%   aideF      : aide pour la série F   (contenu LaTeX)   [optionnel]
%   aideA      : aide pour la série A   (contenu LaTeX)   [optionnel]
%   aideE      : aide pour la série E   (contenu LaTeX)   [optionnel]
%   eleve      : nom de l'élève                           [optionnel]
%   notions    : notions associées (texte libre)          [optionnel]
%
% La fiche utilise les macros Get* pour lire le libellé de l'objectif,
% le critère de réussite (MaitriseS = niveau satisfaisant attendu)
% et le nom de la séquence depuis les bases .dbtex.
% Requiert \seqLoadData préalablement appelé.
%
% Usage minimal :
%   \seqPPRE[
%     niveau=N11, sequence=S03, objectif=02,
%     exosF={1,2,3}, exosA={4,5}, exosE={6},
%     eleve={Alice MARTIN},
%     aideF={Pour l'exercice 1 : dessine d'abord le tableau.},
%     aideA={Pour l'exercice 4 : pense à utiliser le produit en croix.},
%   ]
%
% -----------------------------------------------------------------------
%    \begin{macro}{\seqPPRE}%
%% Options
\define@cmdkey[seq]{PPRE}{objectif}{}%
\define@cmdkey[seq]{PPRE}{exosF}{}%
\define@cmdkey[seq]{PPRE}{exosA}{}%
\define@cmdkey[seq]{PPRE}{exosE}{}%
\define@cmdkey[seq]{PPRE}{aideF}[]{}%
\define@cmdkey[seq]{PPRE}{aideA}[]{}%
\define@cmdkey[seq]{PPRE}{aideE}[]{}%
\define@cmdkey[seq]{PPRE}{eleve}[]{}%
\define@cmdkey[seq]{PPRE}{notions}[]{}%
%
\newcommand{\seqPPRE}[1][]{%
    \setkeys[seq]{PPRE}{aideF,aideA,aideE,eleve,notions,#1}%
    %% ---- Saut de page systématique (une page par objectif)
    \ifdocstd\newpage\fi
    %% ---- Couleurs du thème de la séquence
    \seqSetColors%
    %% ---- En-têtes et pieds de page
    \ifdocstd
        \pagestyle{fancy}\fancyhf{}%
        \lhead{\small PPRE~— \seqNiveauGetNomCourt{\seq@code@niveau}~— \annee}%
        \chead{\small \nomCollege~— \nomClasse}%
        \rhead{\small Objectif~\cmdseq@PPRE@objectif}%
        \lfoot{\small Le \today}%
        \rfoot{\small Page \thepage}%
    \fi
    %% ---- Bandeau titre
    \begin{boiteTitreGen}{%
        \seqSequenceGetNom%
        \hfill
        Objectif~\cmdseq@PPRE@objectif%
    }%
        \begin{minipage}[c][52pt][c]{\linewidth}%
            \vfill
            {\large\bfseries\color{\seq@couleuri@vif}%
                \seqObjectifGetNom{\seq@code@niveau}{\seq@code@sequence}{\cmdseq@PPRE@objectif}}%
            \vfill
            \small
            \ifx\cmdseq@PPRE@eleve\empty\else
                \textbf{Élève~:} \cmdseq@PPRE@eleve\hfill
            \fi
            Niveau actuel~: \fbox{\strut\;F\;}\;\fbox{\strut\;A\;}\;\fbox{\strut\;E\;}%
            \quad (encadrer le niveau noté)%
        \end{minipage}%
    \end{boiteTitreGen}%
    \vspace{4pt}%
    %% ---- Critère de réussite (niveau satisfaisant = niveau attendu)
    \begin{boitePaleNoBreakNoTitle}%
        \small\textit{Pour atteindre le niveau satisfaisant (A)~:}\quad
        \seqObjectifGetMaitriseS{\seq@code@niveau}{\seq@code@sequence}{\cmdseq@PPRE@objectif}%
        \ifx\cmdseq@PPRE@notions\empty\else
            \par\vspace{2pt}%
            \textit{Notions associées~:}\quad \cmdseq@PPRE@notions%
        \fi
    \end{boitePaleNoBreakNoTitle}%
    \vspace{4pt}%
    %% ---- Bloc Série Fondamentale (F)
    \begin{seqPPRE@boite}%
        {Série fondamentale~F — exercices~: \cmdseq@PPRE@exosF}%
        {\seqCouleurSerieF}%
        {\seqCouleurSerieFCadre}%
    \small
    \textit{Ces exercices ont des corrigés en fin de livret.
            Vérifie tes réponses et mesure ce que tu sais déjà faire.}\par
    \ifthenelse{\isempty{\cmdseq@PPRE@aideF}}{%
        \vspace{18pt}%
    }{%
        \vspace{4pt}%
        \begin{boitePaleNoBreakNoTitle}%
            \footnotesize\textbf{Aide~:}\quad \cmdseq@PPRE@aideF%
        \end{boitePaleNoBreakNoTitle}%
    }%
    \end{seqPPRE@boite}%
    \vspace{4pt}%
    %% ---- Bloc Série Avancée (A)
    \begin{seqPPRE@boite}%
        {Série avancée~A — exercices~: \cmdseq@PPRE@exosA}%
        {\seqCouleurSerieA}%
        {\seqCouleurSerieACadre}%
    \small
    \textit{Ces exercices ne sont pas corrigés dans le livret.
            Montre ton travail à l'enseignant.}\par
    \ifthenelse{\isempty{\cmdseq@PPRE@aideA}}{%
        \vspace{18pt}%
    }{%
        \vspace{4pt}%
        \begin{boitePaleNoBreakNoTitle}%
            \footnotesize\textbf{Aide~:}\quad \cmdseq@PPRE@aideA%
        \end{boitePaleNoBreakNoTitle}%
    }%
    \end{seqPPRE@boite}%
    \vspace{4pt}%
    %% ---- Bloc Série Exploration (E)
    \begin{seqPPRE@boite}%
        {Série exploration~E — exercices~: \cmdseq@PPRE@exosE}%
        {\seqCouleurSerieE}%
        {\seqCouleurSerieECadre}%
    \small
    \textit{Ces exercices vont au-delà du niveau attendu.}\par
    \ifthenelse{\isempty{\cmdseq@PPRE@aideE}}{%
        \vspace{18pt}%
    }{%
        \vspace{4pt}%
        \begin{boitePaleNoBreakNoTitle}%
            \footnotesize\textbf{Aide~:}\quad \cmdseq@PPRE@aideE%
        \end{boitePaleNoBreakNoTitle}%
    }%
    \end{seqPPRE@boite}%
}%
%    \end{macro}%
% -------------------------------------------------------------

%% ---- \seqPPRESequence — génère une fiche par objectif de la séquence
%    \begin{macro}{\seqPPRESequence}%
%    \IndexMacro{\seqPPRESequence}%
%
% Génère automatiquement une fiche PPRE pour TOUS les objectifs d'une
% séquence en bouclant sur la base filtrée.
% Les exercices F/A/E doivent être fournis via une macro de dispatch :
%   \cmdExos{codeObjectif}{F|A|E} → numéros des exercices
% Les aides sont optionnelles via une macro de dispatch :
%   \cmdAide{codeObjectif}{F|A|E} → contenu LaTeX (peut être vide)
%
% #1 (optionnel) : options clé=valeur
%   niveau   : code niveau (ex. N11)
%   sequence : code séquence (ex. S03)
%   eleve    : nom de l'élève (optionnel)
% #2 : macro de dispatch des exercices \cmdExos
% #3 : macro de dispatch des aides \cmdAide
%
% Exemple de macros de dispatch :
%   \newcommand{\mesExos}[2]{%
%     \ifthenelse{\equal{#1}{02}}{%
%       \ifthenelse{\equal{#2}{F}}{1,2,3}{%
%       \ifthenelse{\equal{#2}{A}}{4,5}{6}}%
%     }{}% objectifs inconnus : vide
%   }%
%   \newcommand{\mesAides}[2]{%
%     \ifthenelse{\equal{#1}{02}\and\equal{#2}{F}}{%
%       Pour l'exercice 1 : dessine d'abord le triangle.
%     }{}%
%   }%
%   \seqPPRESequence[niveau=N11,sequence=S03,eleve={Alice MARTIN}]
%     {\mesExos}{\mesAides}%
%
% \begin{macrocode}%
\define@cmdkey[seq]{PPRESeq}{eleve}[]{}%
%
\newcommand{\seqPPRESequence}[3][]{%
    \setkeys[seq]{PPRESeq}{eleve,#1}%
    \seq@bases@filter{\seq@code@niveau}{\seq@code@sequence}%
    \DTLforeach*{seq@objectifs@sequence}{%
        \seqPPRE@tmpCode=CodeObjectif%
    }{%
        \seqPPRE@tmpExosF{}%  sera redéfini par #2
        \seqPPRE@tmpExosA{}%
        \seqPPRE@tmpExosE{}%
        \seqPPRE@tmpAideF{}%
        \seqPPRE@tmpAideA{}%
        \seqPPRE@tmpAideE{}%
        % Appel des macros de dispatch
        \edef\seqPPRE@tmpExosF{#2{\seqPPRE@tmpCode}{F}}%
        \edef\seqPPRE@tmpExosA{#2{\seqPPRE@tmpCode}{A}}%
        \edef\seqPPRE@tmpExosE{#2{\seqPPRE@tmpCode}{E}}%
        \def\seqPPRE@tmpAideF{#3{\seqPPRE@tmpCode}{F}}%
        \def\seqPPRE@tmpAideA{#3{\seqPPRE@tmpCode}{A}}%
        \def\seqPPRE@tmpAideE{#3{\seqPPRE@tmpCode}{E}}%
        \seqPPRE[%
            % niveau et sequence lus depuis \seq@code@niveau / \seq@code@sequence%
            %
            objectif = \seqPPRE@tmpCode,%
            exosF    = \seqPPRE@tmpExosF,%
            exosA    = \seqPPRE@tmpExosA,%
            exosE    = \seqPPRE@tmpExosE,%
            aideF    = \seqPPRE@tmpAideF,%
            aideA    = \seqPPRE@tmpAideA,%
            aideE    = \seqPPRE@tmpAideE,%
            eleve    = \cmdseq@PPRESeq@eleve,%
        ]%
    }%
}%
% \end{macrocode}%
%    \end{macro}%

% Variables temporaires internes pour \seqPPRESequence
\newcommand{\seqPPRE@tmpCode}{}%
\newcommand{\seqPPRE@tmpExosF}{}%
\newcommand{\seqPPRE@tmpExosA}{}%
\newcommand{\seqPPRE@tmpExosE}{}%
\newcommand{\seqPPRE@tmpAideF}{}%
\newcommand{\seqPPRE@tmpAideA}{}%
\newcommand{\seqPPRE@tmpAideE}{}%


%% ------------------------------------------------------------------
%% ---- Livret de séquence ------------------------------------------
%% ------------------------------------------------------------------

%% ---- Titre de section --------------------------------------------
%    \begin{macro}{\seqTitreSection}%
\newcommand{\seqTitreSection}[1]{%
	\ifdocstd
		\newpage
		\newgeometry{left=35pt,right=35pt,top=65pt,bottom=70pt}%
	\fi
	\begin{boiteTitreSection}%
		\begin{minipage}[c][20pt][c]{\linewidth}%
			\centering \Large \textbf{#1}%
		\end{minipage}%
	\end{boiteTitreSection}%
	\par\vspace{10pt}%
}%
%    \end{macro}%
% -------------------------------------------------------------

%    \begin{macro}{\seqTitreSectionCorrige}%
\newcommand{\seqTitreSectionCorrige}[1]{%
	\par\vspace{10pt}%
	\begin{boiteTitreSection}%
		\begin{minipage}[c][20pt][c]{\linewidth}%
			\centering \Large \textbf{#1}%
		\end{minipage}%
	\end{boiteTitreSection}%
	\par\vspace{5pt}%
}%
%    \end{macro}%
% -------------------------------------------------------------

%% ---- Pour chaque annexe -----------------------------------------
%    \begin{macro}{\seqAnnexe}%
%% Ecrit le corrigé de l'exercice dans un fichier séparé
\newcommand{\seqAnnexe}[1]{%
	\refstepcounter{AnnexeNum}\label{annexe:\theAnnexeNum}\theAnnexeNum\kern-0.5em%
	\immediate\write\annfile{%
		\unexpanded{\par\nopagebreak\vspace{5pt}\nopagebreak%
			\begin{boiteTitreExo}}%
				Annexe\noexpand\;\theAnnexeNum
				\unexpanded{\end{boiteTitreExo}%
			\par\nopagebreak\vspace{5pt}\nopagebreak}%
		\noexpand\input{Annexes/a\theAnnexeNum.tex}%
	}%
	\scantokens{\begin{tcbverbatimwrite}%
			{Annexes/a\theAnnexeNum.tex} #1
		\end{tcbverbatimwrite}%
	}%
}%
%    \end{macro}%
% -------------------------------------------------------------

%% ---- Regroupe les annexes des exercices ----------------------
%    \begin{macro}{\seqAfficheAnnexes}%
\newcommand{\seqAfficheAnnexes}{ \immediate\closeout\annfile \ifnumcomp{\theAnnexeNum}{=}{0}{}{\newpage \seqTitreSectionCorrige {Annexes}
\newpage}}%
%    \end{macro}%

%% ---- v0.11.5 — Macros liées aux exercices déplacées dans
%%                seqenseigne-core-exos.dtx :
%%                  \seqAfficheCorriges, \seqCorrige, \seqCorrigesExos,
%%                  \seqInitCorriges, \seqAfficheCorrigesBilan (reste ici),
%%                  \seq@ecrit@corrige@serieExos, \seq@ecrit@corrige@Exo,
%%                  seqSerieExos, seqExercice.
%%                Le module core-exos est chargé par seqenseigne.sty
%%                APRÈS ce module core, donc \seqInitCorriges utilisée par
%%                \seqTitreLivret plus bas est bien définie au moment où
%%                l'utilisateur invoque \seqTitreLivret dans son document.
% -------------------------------------------------------------

% ---- Gestion des compteurs
\newcommand\seqCreeCompteurs{%
	\newcounter{NotionNum}%
	\newcounter{DefinitionNum}%
	\newcounter{ProprieteNum}%
	\newcounter{MethodeNum}%
	\newcounter{SerieExosNum}%
	\newcounter{ExoNum}%
	\newcounter{AnnexeNum}%
	%	\newcounter{SousSection}%
	%% --- Ancres hyperref uniques : theH<compteur> inclut le numéro de série
	%% Sans cela, hyperref génère "ExoNum.1" à chaque série → warnings
	%% "destination with the same identifier"
	\renewcommand{\theHExoNum}{\theSerieExosNum-\theExoNum}%
	\renewcommand{\theHNotionNum}{\theSerieExosNum-\theNotionNum}%
	\renewcommand{\theHMethodeNum}{\theSerieExosNum-\theMethodeNum}%
	\renewcommand{\theHDefinitionNum}{\theSerieExosNum-\theDefinitionNum}%
	\renewcommand{\theHProprieteNum}{\theSerieExosNum-\theProprieteNum}%
}%

\newcommand\seqRAZCompteurs{%
	\setcounter{NotionNum}{0}%
	\setcounter{DefinitionNum}{0}%
	\setcounter{ProprieteNum}{0}%
	\setcounter{MethodeNum}{0}%
	\setcounter{SerieExosNum}{0}%
	\setcounter{ExoNum}{0}%
	\setcounter{AnnexeNum}{0}%
	%	\setcounter{SousSection}{0}%
}%

%% ---- Initialisation des annexes ---------------------------------
\newcommand{\seqInitAnnexes}{%
	%% ---- Création du fichier pour les annexes
	\newwrite\annfile
	\immediate\openout\annfile=AnnList.tex
	\immediate\write\annfile{\noexpand\newpage\noexpand\seqTitreSectionCorrige{Annexes}}%
}%

%% ---- Titre du livret de séquence ---------------------------------
%    \begin{macro}{\seqTitreLivret}%
%% ---- Options 
%    \begin{macrocode}%
\define@cmdkey[seq]{titreLivret}{corriges}[non]{}%
\newcommand{\seqTitreLivret}[1][]{%
	%% ---- Options
	\setkeys[seq]{titreLivret}{corriges,#1}%
	%% ---- Initialisation des corrigés et des annexes
	\seqInitCorriges
	\seqInitAnnexes
	%% ---- Corrigés ?
	\ifthenelse{\equal{\cmdseq@titreLivret@corriges}{oui}}%
		{\renewcommand{\seqCorrigesExos}{oui}}%
		{\renewcommand{\seqCorrigesExos}{non}}%
	%% ---- Entêtes et pieds de page
	\ifdocstd
		\pagestyle{fancy}%
		\fancyhf{}%
		\lfoot{}%
		\rfoot{Page \thepage/\pageref{LastPage}}%
		\lfoot{\seqSequenceGetNom}%
	\fi
	%% ---- Couleurs
	\seqSetColors
	%% ---- Compteurs
	\seqCreeCompteurs
	%% ---- Boîte de titre
	\begin{boiteTitreGen}{Séquence \seqSequenceGetNumero}%
		\begin{minipage}[c][100pt][c]{\linewidth}%
			\vfill \textit{Livret de séquence} \hfill
			\vfill \centering {\Huge \textbf{\seqSequenceGetNom}}%
			\vfill Version de \datemoisannee\today \hfill \textit{Classe de \seqNiveauGetNomCourt{\seq@code@niveau}}%
		\end{minipage}%
	\end{boiteTitreGen}%
}%
%    \end{macro}%
% -------------------------------------------------------------

%% ---- Titre du livret de séquence (mode auto) ---------------------
%    \begin{macro}{\seqTitrelivretauto}%
%    \begin{macrocode}%
\define@cmdkey[seq]{titreLivretAuto}{corriges}[non]{}%
\newcommand{\seqTitreLivretAuto}[1][]{%
	%% ---- Options
	\setkeys[seq]{titreLivretAuto}{corriges,#1}%
	\seqTitreLivret[corriges=\cmdseq@titreLivretAuto@corriges]
}%
%    \end{macro}%

%% v0.13.5.2 — Code des évaluations (Bilan modèle 2021 + Eval modèle 2019) déplacé.
%%   - Bilan 2021    → seqenseigne-legacy.dtx
%%   - Eval 2019     → seqenseigne-core-eval.dtx
%%   - QCM (env.)    → seqenseigne-core-exos.dtx
%% Voir aussi le bloc final (compteurs/writes) en fin de fichier.

%% ------------------------------------------------------------------
%% ---- Bilan et Évaluation -----------------------------------------
%% ------------------------------------------------------------------
%% v0.13.5.2 étape 1 — Déplacement du code des évaluations sommatives :
%%   - Modèle Bilan 2021 → seqenseigne-legacy.dtx
%%       (\seqBilan, \seqBilanItemAlea, \seqBilanCorrigeExo, ...)
%%   - Modèle Eval 2019  → seqenseigne-core-eval.dtx
%%       (\titreEval renommé \seqTitreEval en étape 2,
%%        \corrigeExoEval renommé \seqEvalCorrigeExo en étape 2,
%%        \seqEvalBaremeItem, \seqEvalAfficheCorriges, ...)
%%   - Environnement qcm → seqenseigne-core-exos.dtx (renommé seqQcm
%%        en étape 2, refactorisé en environnement enfant)
%%
%% v0.13.5.2 étape 2 — Centralisation des compteurs (cf. fin de fichier) :
%%   tous les \newcounter sont désormais dans core, indépendamment du
%%   module qui les utilise.
%%
%% Le \newwrite\corrfileeval, partagé entre eval (core-eval) et
%% bilan (legacy), reste dans core (squelette commun, cf. fin de
%% fichier).
% -------------------------------------------------------------



% =============================================================
% Environnements structurants
% =============================================================
% Ici tu colles les environnements qui organisent le contenu
% pédagogique (par ex. les séries d'exercices, les blocs de cours,
% les plans de travail…), quand leur *logique* n'est pas purement visuelle.
%
% Exemple schématique :
%
% \begin{environment}{seqSerieExos}%
% \IndexEnv{seqSerieExos}%
% Environnement listant une série d'exercices.
% \begin{macrocode}%
% \newenvironment{seqSerieExos}[1][]{%
%   % code début d'environnement
% }{%
%   % code fin d'environnement
% }%
% \end{macrocode}%
% \end{environment}%
%
%% ------------------------------------------------------------------

%% ------------------------------------------------------------------
%% ---- Plan de travail ---------------------------------------------
%% ------------------------------------------------------------------

%    \begin{environment}{titrePlan}%
%% Options
\define@cmdkey[seq]{titrePlan}{titreSeq}{}%
\define@cmdkey[seq]{titrePlan}{numSeq}{}%
\define@cmdkey[seq]{titrePlan}{classe}{}%
\define@cmdkey[seq]{titrePlan}{etab}{}%
\define@cmdkey[seq]{titrePlan}{annee}{}%
\define@cmdkey[seq]{titrePlan}{enseignant}{}%
\define@cmdkey[seq]{titrePlan}{theme}[gris]{}%
%% ---- Titre du plan de travail ------------------------------------
\newenvironment{titrePlan}[1][]{%
	%% ---- Options
	\setkeys[seq]{titrePlan}{theme,#1}%
	%% ---- Entêtes et pieds de page
	\ifdocstd
		\pagestyle{fancy}%
		\fancyhf{}%
		\lfoot{Le \today}%
		\rfoot{Page \thepage/\pageref{LastPage}}%
	\fi
	%% ---- Couleurs
	\seq@setColors{\cmdseq@titrePlan@theme}%
	%% ---- Boîte de titre
	\begin{boiteTitreGen}{Séquence \cmdseq@titrePlan@numSeq}%
		\begin{minipage}[c][70pt][c]{\linewidth}%
			\vfill \textit{Plan de travail}%
			\ifx\cmdseq@titrePlan@annee\undefined\else
				\hfill \textit{Année \cmdseq@titrePlan@annee}%
			\fi
			\vfill \centering {\Huge \textbf{\cmdseq@titrePlan@titreSeq}}%
			\vfill
			\ifx\cmdseq@titrePlan@enseignant\undefined\else
				\textit{Enseignant : \cmdseq@titrePlan@enseignant}%
			\fi
			\hfill \textit{\cmdseq@titrePlan@etab - \cmdseq@titrePlan@classe}%
		\end{minipage}%
	\end{boiteTitreGen}%
	\xltabular{\linewidth}{*{2}{p{0.45\linewidth}}}%
}{%
	\endxltabular
}%
%    \end{environment}%

%    \begin{environment}{blocReperesNiveaux}%
\newenvironment{blocReperesNiveaux}[1]{%
	\begin{boiteVariableWidthHeightTitre}{\large \textbf{#1}}{\linewidth}{170pt}%
		\begin{dinglist}{224}%
			}{%
		\end{dinglist}%
	\end{boiteVariableWidthHeightTitre}%
}%
%    \end{environment}%


%% ------------------------------------------------------------------
%% ---- Livret de séquence ------------------------------------------
%% ------------------------------------------------------------------


%% ---- Objectifs de niveau 1 et 2 ----------------------------------
%    \begin{environment}{seqObjectifs}%
\newenvironment{seqObjectifs}[1][
	Les capacités suivantes sont les objectifs de la séquence :
	ce sont elles qui seront évaluées en fin de séquence.
]{%
	\par\vspace{20pt}%
	\begin{boitePaleNoBreak}{Objectifs de fin de séquence}%
		#1
		\tcblower
		\begin{minipage}{\linewidth}%
			\begin{multicols}{2} \begin{dinglist}{224}%
					}{%
				\end{dinglist} \end{multicols}%
		\end{minipage}%
	\end{boitePaleNoBreak}%
	\par\vspace{10pt}%
}%
%    \end{environment}%

%    \begin{environment}{seqObjectifsN2}%
\newenvironment{seqObjectifsN2}{%
	\begin{dinglist}{118}%
		}{%
	\end{dinglist}%
}%
%    \end{environment}%

%% ---- Carte mentale des connaissances -----------------------------
%    \begin{environment}{seqCarteMentale}%
\newenvironment{seqCarteMentale}[1][La carte mentale suivante
	représente de manière très synthétique l'ensemble
	des connaissances de la séquence.
]{%
	\par\vspace{20pt}%
	\begin{boitePaleNoBreak}{Carte mentale}%
		#1
		\tcblower
		}{%
	\end{boitePaleNoBreak}%
	\par\vspace{10pt}%
}%
%    \end{environment}%

%% ---- Notions ---------------------------------------------------
%    \begin{environment}{notion}%
\newenvironment{seqNotion}[2]{%
	\refstepcounter{NotionNum}\label{notion:\theNotionNum}%
	\setlength{\columnsep}{10pt}%
	\par\vspace{5pt}%
	\begin{boitePaleBreakFitHeight}{%
			Notion~\theNotionNum
			\ifthenelse{\isempty{#1}}{}%
			{ :\quad \begin{minipage}{0.5\linewidth}#1\end{minipage}}%
		}%
		#2
		\tcblower
		}{%
	\end{boitePaleBreakFitHeight}%
	\par\vspace{5pt}%
}%
%    \end{environment}%

%% ---- Définitions -----------------------------------------------
%    \begin{environment}{definition}%
\newenvironment{seqDefinition}[2]{%
	\refstepcounter{DefinitionNum}\label{def:\theDefinitionNum}%
	\setlength{\columnsep}{10pt}%
	\par\vspace{5pt}%
	\begin{boitePaleBreakFitHeight}{%
			Définition~\theDefinitionNum
			\ifthenelse{\isempty{#1}}{}%
			{ :\quad \begin{minipage}{0.5\linewidth}#1\end{minipage}}%
		}%
		#2
		\tcblower
		}{%
	\end{boitePaleBreakFitHeight}%
	\par\vspace{5pt}%
}%
%    \end{environment}%

%% ---- Propriétés -----------------------------------------------
%    \begin{environment}{propriete}%
\newenvironment{seqPropriete}[2]{%
	\refstepcounter{ProprieteNum}\label{prop:\theProprieteNum}%
	\setlength{\columnsep}{10pt}%
	\par\vspace{5pt}%
	\begin{boitePaleBreakFitHeight}{%
			Propriété~\theProprieteNum
			\ifthenelse{\isempty{#1}}{}%
			{ :\quad \begin{minipage}{0.5\linewidth}#1\end{minipage}}%
		}%
		#2
		\tcblower
		}{%
	\end{boitePaleBreakFitHeight}%
	\par\vspace{5pt}%
}%
%    \end{environment}%

%% ---- Méthodes -----------------------------------------------
%    \begin{environment}{methode}%
\newenvironment{seqMethode}[2]{%
	\refstepcounter{MethodeNum}\label{methode:\theMethodeNum}%
	\setlength{\columnsep}{10pt}%
	\par\vspace{5pt}%
	\begin{boitePaleBreakFitHeight}{%
			Méthode~\theMethodeNum
			\ifthenelse{\isempty{#1}}{}%
			{ :\quad \begin{minipage}{0.5\linewidth}#1\end{minipage}}%
		}%
		#2
		\tcblower
		}{%
	\end{boitePaleBreakFitHeight}%
	\par\vspace{5pt}%
}%
%    \end{environment}%

%    \begin{environment}{seqSerieExos}%
%% ---- v0.11.5 — seqSerieExos et seqExercice déplacés dans
%%                seqenseigne-core-exos.dtx, ainsi que les helpers
%%                \seq@ecrit@corrige@serieExos et \seq@ecrit@corrige@Exo
%%                qui les accompagnent.
% -------------------------------------------------------------

%% ------------------------------------------------------------------
%% v0.13.5.2 — Environnements eval/bilan déplacés (cf. note plus haut).




% =============================================================
% Macros internes \seq@...
% =============================================================
% Les helpers internes non destinés à l'utilisateur.
% Garde leur doc interne, leurs commentaires, etc.
% -------------------------------------------------------------
% \begin{macro}{\seq@TrucInterne}%
% Macro interne.
% \begin{macrocode}%
% \newcommand{\seq@TrucInterne}[2]{...}%
% \end{macrocode}%
% \end{macro}%

% \begin{macro}{\seq@ecrit@corrige@serieExos}%
% \begin{macro}{\seq@ecrit@corrige@Exo}%
%% v0.11.5 — Helpers déplacés dans seqenseigne-core-exos.dtx, avec
%% au passage le patch D4 : \noexpand\newpage → \noexpand\clearpage
%% dans \seq@ecrit@corrige@serieExos pour forcer la sortie des flottants
%% en attente entre deux séries de corrigés.
% \end{macro}%
% \end{macro}%
% -------------------------------------------------------------

%% v0.13.5.2 — \seq@Eval@QcmItemNum déplacé dans seqenseigne-core-eval.dtx,
%%             \repCodes déplacé dans seqenseigne-core-exos.dtx.



% -------------------------------------------------------------
% Hooks divers, patchs, compat low-level, mesures, longueurs...
% (toutes les \newlength, \setlength, \newcounter, etc.
% qui ne sont pas purement de la déco)
% -------------------------------------------------------------

%% ------------------------------------------------------------------
%% ---- Compteurs et writes centralisés -----------------------------
%% ------------------------------------------------------------------
%% v0.13.5.2 étape 2 : centralisation de tous les compteurs dans core,
%% pour visibilité globale et pour éviter les dépendances croisées de
%% \newcounter entre modules.

%% ---- Compteurs pour les QCM (utilisés par core-exos)
%% Note v0.13.5.2 étape 2 : le compteur seqQcmItemNum a été supprimé.
%% La numérotation des items du QCM se fait désormais via \therownum
%% (compteur fourni par tabularray, robuste face aux passes multiples
%% du moteur), avec décalage de -2 pour les 2 lignes d'en-tête.
%% Les helpers de construction de l'en-tête A B C... restent.

%% Helpers internes pour la construction de l'en-tête de colonnes
%% (A, B, C...) du QCM. seq@qcm@iter sert de variable de boucle,
%% seq@qcm@alphcode sert à générer le code ASCII des lettres.
%% Renommés depuis i et AlphCode (v0.13.5.2 étape 2) pour préfixage
%% cohérent et pour éviter toute collision avec un éventuel \i ou
%% AlphCode utilisateur.
\newcounter{seq@qcm@iter}%
\newcounter{seq@qcm@alphcode}%

%% ---- Compteur pour le modèle Bilan 2021 (utilisé par legacy)
\newcounter{seqBilanExoNum}%

%% ---- Compteur pour les flashcards (utilisé par flashcard)
\newcounter{flashcardnb}%

%% ------------------------------------------------------------------
%% ---- Fractions avec strut ----------------------------------------
%% ------------------------------------------------------------------
%    \begin{macro}{\seqFrac}%
%    \IndexMacro{\seqFrac}%
%    Fraction avec struts pour un espacement vertical uniforme
%    #1 : numérateur
%    #2 : dénominateur
%    \begin{macrocode}%
\newcommand{\seqFrac}[2]{\dfrac{\strut{} #1}{\strut{} #2}}%
%    \end{macrocode}%
%    \end{macro}%

%% ------------------------------------------------------------------
%% ---- Dénominateur de fraction vers texte -------------------------
%% ------------------------------------------------------------------
\newcommand\seqFracToText[1]{%
	\ifnum#1=2{demi\;}\else
		\ifnum#1=3{tiers\;}\else
			\ifnum#1=4{quart\;}\else
				\ifnum#1=5{cinquième\;}\else
					\ifnum#1=6{sixième\;}\else
						\ifnum#1=7{septième\;}\else
							\ifnum#1=8{huitième\;}\else
								\ifnum#1=9{neuvième\;}\else
									\ifnum#1=10{dixième\;}\else
										{inconnu\;}%
									\fi
								\fi
							\fi
						\fi
					\fi
				\fi
			\fi
		\fi
	\fi
}%

%% ------------------------------------------------------------------
%% ---- Nombre vers texte -------------------------------------------
%% ------------------------------------------------------------------
\newcommand\seqNbToText[1]{%
	\ifnum#1=2{deux\;}\else
		\ifnum#1=3{trois\;}\else
			\ifnum#1=4{quatre\;}\else
				\ifnum#1=5{cinq\;}\else
					\ifnum#1=6{six\;}\else
						{inconnu\;}%
					\fi
				\fi
			\fi
		\fi
	\fi
}%

%% ------------------------------------------------------------------
%% ---- Macros de convenance pour documents de suivi ---------------
%% ------------------------------------------------------------------

% \begin{macro}{\seqTitreSuivi}%
% \IndexMacro{\seqTitreSuivi}%
% Configure les en-têtes/pieds de page pour documents de suivi
% #1 : classe
% #2 : titre
% #3 : année
% \begin{macrocode}%
\newcommand{\seqTitreSuivi}[3]{%
    \pagestyle{fancy}%
    \fancyhf{}%
    \lhead{#1}%
    \chead{#2}%
    \rhead{#3}%
    \lfoot{Le \today}%
    \rfoot{Page \thepage/\pageref{LastPage}}%
    \setlength{\footskip}{20.0pt}%
}%
% \end{macrocode}%
% \end{macro}%

% \begin{macro}{\seqSuiviSequence}%
% \IndexMacro{\seqSuiviSequence}%
% Génère le suivi d'une séquence unique
% #1 : code niveau
% #2 : code séquence
% \begin{macrocode}%
\newcommand{\seqSuiviSequence}[2]{%
    \seqSetCodeNiveau{#1}%
    \seqSetCodeSequence{#2}%
    \chead{Suivi périodique: #2~-~\seqSequenceGetNom}%
    \seqSuiviAnnuelParSequence[sequence=#2,niveau=#1]\newpage
}%
% \end{macrocode}%
% \end{macro}%

% Fin du module core
%</package>
